\section{単位系・換算}

\subsection{平成23年度 問6}
\label{q:h23-reg-applied-06}

\noindent\textbf{原典：} sources/H23/17\_applied\_science.pdf、PDF物理3頁。

\textbf{公式正答：} 3（single）。

\par\textbf{検討未確定：} 原典の導入は6基本単位だけを列挙し、光度のcdを欠いてその他全部を組立できるとする。SIは出題時も7基本単位である。選択肢内の指数誤りは特定できるが導入の全称を正しい記述として校正完了にしない。

\paragraph{問題の要旨} 放射能・力・磁束密度・仕事率・抵抗を基本単位へ分解し、指数の誤りを選ぶ。

\Needspace{4\baselineskip}\paragraph{確認した条件・記述（独自要約）}\mbox{}\par

\noindent\textbf{方向：} 誤っているもの。\par

\noindent\textbf{対象：} SIの定義・換算。定義の時点と指数を全肢画像で照合。\par

\paragraph{問われる核心} 物理量の定義から組立単位を導く

\paragraph{必要知識} F=\allowbreak{}ma、p=\allowbreak{}F/S、仕事Fℓ、P=\allowbreak{}E/t、Q=\allowbreak{}It、V=\allowbreak{}E/Q、R=\allowbreak{}V/I、磁気力IℓB。

\paragraph{考え方・正答を選ぶ理由（査読待ち初稿）} 3が指数を取り違えている。\(F=I\ell B\)より\([B]=\mathrm{N/(A\,m)=kg\,s^{-2}A^{-1}}\)である。長さを一つ消し電流指数を-1にする必要がある。Bq/N/W/Ωは各物理定義から導ける。導入のcd欠落は別の疑義として保持する。

\Needspace{4\baselineskip}\paragraph{選択肢の検討（査読待ち初稿）}\mbox{}\par

\noindent\textbf{1：適切} \(\mathrm{Bq=s^{-1}}\)。放射能は単位時間の壊変数なのでBqはsの-1乗。周期周波数と同じ次元でも意味は区別する。\par

\noindent\textbf{2：適切} \(\mathrm{N=m\,kg\,s^{-2}}\)。力F=\allowbreak{}maからN=\(\mathrm{kg\,m\,s^{-2}}\)である。\par

\noindent\textbf{3：不適切} \(\mathrm{T=m\,kg\,s^{-2}A^{-2}}\)。電流に働く磁気力F=\allowbreak{}IℓBからT=\allowbreak{}N/(A m)=\allowbreak{}kg s\textasciicircum{}-2 A\textasciicircum{}-1。原典は不要なmとAの指数-2を含む。\par

\noindent\textbf{4：適切} \(\mathrm{W=m^2kg\,s^{-3}}\)。仕事率はエネルギーを時間で割りW=\(\mathrm{kg\,m^2s^{-3}}\)である。\par

\noindent\textbf{5：適切} \(\mathrm{\Omega=m^2kg\,s^{-3}A^{-2}}\)。抵抗R=\allowbreak{}V/IからΩ=\(\mathrm{kg\,m^2s^{-3}A^{-2}}\)である。\par

\paragraph{再利用する考え方} 定義の基準日を確定し、物理量の関係から指数と係数を導く。量子現象では運動量・エネルギー保存と観測対象を区別する。

\paragraph{共通理論} \hyperref[theory:units:derived]{第\ref*{theory:units:derived}節}, \hyperref[theory:units:electrical]{第\ref*{theory:units:electrical}節}, \hyperref[theory:units:activity]{第\ref*{theory:units:activity}節}

\par 参考：\cite{si_nist811,si_brochure9}。

\paragraph{関連問題} \hyperref[q:h27-reg-applied-04]{H27-REG-APPLIED-04}, \hyperref[q:r05-reg-applied-13]{R05-REG-APPLIED-13}

\par 難易度 1（校準初稿）、優先度未評価。

\subsection{平成25年度 問4}
\label{q:h25-reg-applied-04}

\noindent\textbf{原典：} sources/H25/17\_applied\_science.pdf、PDF物理3頁。

\textbf{公式正答：} 5（single）。

\paragraph{問題の要旨} 周波数・力・圧力・仕事・電圧の組立単位を比較する。

\Needspace{4\baselineskip}\paragraph{確認した条件・記述（独自要約）}\mbox{}\par

\noindent\textbf{方向：} 誤っているもの。\par

\noindent\textbf{対象：} SIの定義・換算。定義の時点と指数を全肢画像で照合。\par

\paragraph{問われる核心} 電気量を力学的エネルギーへ接続する

\paragraph{必要知識} Hz=\allowbreak{}1/T、N=\allowbreak{}ma、Pa=\allowbreak{}N/m\textsuperscript{2}、J=\allowbreak{}Nm、C=\allowbreak{}As、V=\allowbreak{}J/C。

\paragraph{考え方・正答を選ぶ理由（査読待ち初稿）} 5の長さ指数が不足している。電圧は1Cを移す仕事で\(V=J/C\)、\([V]=\mathrm{kg\,m^2s^{-3}A^{-1}}\)となる。OCRのma等を文字通り使わず原典上付きと全7基本単位を確認した。

\Needspace{4\baselineskip}\paragraph{選択肢の検討（査読待ち初稿）}\mbox{}\par

\noindent\textbf{1：適切} \(\mathrm{Hz=s^{-1}}\)。周期Tの逆数でHz=\allowbreak{}s\textasciicircum{}-1である。\par

\noindent\textbf{2：適切} \(\mathrm{N=m\,kg\,s^{-2}}\)。力F=\allowbreak{}maからN=\(\mathrm{kg\,m\,s^{-2}}\)である。\par

\noindent\textbf{3：適切} \(\mathrm{Pa=m^{-1}kg\,s^{-2}}\)。圧力p=\allowbreak{}F/SからPa=\(\mathrm{kg\,m^{-1}s^{-2}}\)である。\par

\noindent\textbf{4：適切} \(\mathrm{J=m^2kg\,s^{-2}}\)。仕事は力と移動距離の積でJ=\(\mathrm{kg\,m^2s^{-2}}\)である。\par

\noindent\textbf{5：不適切} \(\mathrm{V=m\,kg\,s^{-3}A^{-1}}\)。電位差は仕事/電荷からV=\allowbreak{}J/C=\(\mathrm{kg\,m^2s^{-3}A^{-1}}\)である。原典はmの指数が1であり2ではない。\par

\paragraph{再利用する考え方} 定義の基準日を確定し、物理量の関係から指数と係数を導く。量子現象では運動量・エネルギー保存と観測対象を区別する。

\paragraph{共通理論} \hyperref[theory:units:derived]{第\ref*{theory:units:derived}節}, \hyperref[theory:units:electrical]{第\ref*{theory:units:electrical}節}

\par 参考：\cite{si_nist811,si_brochure9}。

\paragraph{関連問題} \hyperref[q:h27-reg-applied-04]{H27-REG-APPLIED-04}

\par 難易度 1（校準初稿）、優先度未評価。

\subsection{平成27年度 問4}
\label{q:h27-reg-applied-04}

\noindent\textbf{原典：} sources/H27/17\_applied\_science.pdf、PDF物理2頁。

\textbf{公式正答：} 4（single）。

\par\textbf{検討未確定：} 原典の導入は6基本単位だけを列挙し、光度のcdを欠いてその他全部を組立できるとする。SIは出題時も7基本単位である。選択肢内の指数誤りは特定できるが導入の全称を正しい記述として校正完了にしない。

\paragraph{問題の要旨} 電荷・力・圧力・仕事率・抵抗の基本単位指数を確認する。

\Needspace{4\baselineskip}\paragraph{確認した条件・記述（独自要約）}\mbox{}\par

\noindent\textbf{方向：} 最も不適切なもの。\par

\noindent\textbf{対象：} SIの定義・換算。定義の時点と指数を全肢画像で照合。\par

\paragraph{問われる核心} 仕事率の長さ指数と電気量の結合

\paragraph{必要知識} F=\allowbreak{}ma、仕事Fℓ、P=仕事/t、Q=\allowbreak{}It、R=\allowbreak{}V/I。

\paragraph{考え方・正答を選ぶ理由（査読待ち初稿）} 4が仕事率を誤っている。\(P=F\ell/t\)から\([P]=\mathrm{kg\,m^2s^{-3}}\)でありm\textsuperscript{2}が必要である。電荷と抵抗も電流と時間・仕事から同じ体系で導く。導入はcdを欠くため疑義を別に残す。

\Needspace{4\baselineskip}\paragraph{選択肢の検討（査読待ち初稿）}\mbox{}\par

\noindent\textbf{1：適切} \(\mathrm{C=s\,A}\)。電荷Q=\allowbreak{}ItよりC=\allowbreak{}Asである。\par

\noindent\textbf{2：適切} \(\mathrm{N=m\,kg\,s^{-2}}\)。力F=\allowbreak{}maからN=\(\mathrm{kg\,m\,s^{-2}}\)である。\par

\noindent\textbf{3：適切} \(\mathrm{Pa=m^{-1}kg\,s^{-2}}\)。圧力p=\allowbreak{}F/SからPa=\(\mathrm{kg\,m^{-1}s^{-2}}\)である。\par

\noindent\textbf{4：不適切} \(\mathrm{W=m\,kg\,s^{-3}}\)。仕事率はエネルギーを時間で割りW=\(\mathrm{kg\,m^2s^{-3}}\)である。原典は長さ指数1であり、仕事の移動距離が一つ不足する。\par

\noindent\textbf{5：適切} \(\mathrm{\Omega=m^2kg\,s^{-3}A^{-2}}\)。抵抗R=\allowbreak{}V/IからΩ=\(\mathrm{kg\,m^2s^{-3}A^{-2}}\)である。\par

\paragraph{再利用する考え方} 定義の基準日を確定し、物理量の関係から指数と係数を導く。量子現象では運動量・エネルギー保存と観測対象を区別する。

\paragraph{共通理論} \hyperref[theory:units:derived]{第\ref*{theory:units:derived}節}, \hyperref[theory:units:electrical]{第\ref*{theory:units:electrical}節}

\par 参考：\cite{si_nist811,si_brochure9}。

\paragraph{関連問題} \hyperref[q:h23-reg-applied-06]{H23-REG-APPLIED-06}, \hyperref[q:h25-reg-applied-04]{H25-REG-APPLIED-04}

\par 難易度 1（校準初稿）、優先度未評価。

\subsection{令和5年度 問13}
\label{q:r05-reg-applied-13}

\noindent\textbf{原典：} sources/R05/17\_applied\_science.pdf、PDF物理5頁。

\textbf{公式正答：} 5（single）。

\paragraph{問題の要旨} SI組立単位と動粘度の非SI単位の換算を比較する。

\Needspace{4\baselineskip}\paragraph{確認した条件・記述（独自要約）}\mbox{}\par

\noindent\textbf{方向：} 最も不適切なもの。\par

\noindent\textbf{対象：} SIの定義・換算。定義の時点と指数を全肢画像で照合。\par

\paragraph{問われる核心} 組立単位の指数と非SI換算係数

\paragraph{必要知識} C=\allowbreak{}As、J=\allowbreak{}Nm、Pa=\allowbreak{}N/m\textsuperscript{2}、Ω=V/A、動粘度St=\allowbreak{}cm\textsuperscript{2}/s。

\paragraph{考え方・正答を選ぶ理由（査読待ち初稿）} 5の電流指数が誤りである。\(R=V/I=(J/Q)/I\)、\(Q=It\)より\([R]=\mathrm{kg\,m^2s^{-3}A^{-2}}\)。4は非SI単位StをSIへ正しく換算した記述であり、「SIでないから換算も誤り」としない。

\Needspace{4\baselineskip}\paragraph{選択肢の検討（査読待ち初稿）}\mbox{}\par

\noindent\textbf{1：適切} \(1\mathrm{C}=1\mathrm{A\,s}\)。電荷Q=\allowbreak{}ItよりC=\allowbreak{}Asである。\par

\noindent\textbf{2：適切} \(1\mathrm{J}=1\mathrm{kg\,m^2s^{-2}}\)。仕事は力と移動距離の積でJ=\(\mathrm{kg\,m^2s^{-2}}\)である。\par

\noindent\textbf{3：適切} \(1\mathrm{Pa}=1\mathrm{kg\,m^{-1}s^{-2}}\)。圧力p=\allowbreak{}F/SからPa=\(\mathrm{kg\,m^{-1}s^{-2}}\)である。\par

\noindent\textbf{4：適切} \(1\mathrm{St}=10^{-4}\mathrm{m^2s^{-1}}\)。Stは動粘度cm\textsuperscript{2}/s。cm=\allowbreak{}10\textasciicircum{}-2mを二乗するので10\textasciicircum{}-4m\textsuperscript{2}/sである。SI基本単位そのものとは呼ばない。\par

\noindent\textbf{5：不適切} \(1\mathrm{\Omega}=1\mathrm{m^2kg\,s^{-3}A^{-1}}\)。抵抗R=\allowbreak{}V/IからΩ=\(\mathrm{kg\,m^2s^{-3}A^{-2}}\)である。原典はA\textasciicircum{}-1であり抵抗に必要なA\textasciicircum{}-2を欠く。\par

\paragraph{再利用する考え方} 定義の基準日を確定し、物理量の関係から指数と係数を導く。量子現象では運動量・エネルギー保存と観測対象を区別する。

\paragraph{共通理論} \hyperref[theory:units:derived]{第\ref*{theory:units:derived}節}, \hyperref[theory:units:electrical]{第\ref*{theory:units:electrical}節}, \hyperref[theory:units:conversion]{第\ref*{theory:units:conversion}節}

\par 参考：\cite{si_nist811}。

\paragraph{関連問題} \hyperref[q:h23-reg-applied-06]{H23-REG-APPLIED-06}

\par 難易度 1（校準初稿）、優先度未評価。

\subsection{令和6年度 問16}
\label{q:r06-reg-applied-16}

\noindent\textbf{原典：} sources/R06/17\_applied\_science.pdf、PDF物理5頁。

\textbf{公式正答：} 3（single）。

\paragraph{問題の要旨} 圧力の慣用単位をPaへ換算し、最も遠い係数を選ぶ。

\Needspace{4\baselineskip}\paragraph{確認した条件・記述（独自要約）}\mbox{}\par

\noindent\textbf{方向：} 圧力換算として最も遠い値。\par

\noindent\textbf{対象：} SIの定義・換算。定義の時点と指数を全肢画像で照合。\par

\paragraph{問われる核心} 力単位と面積換算の二乗

\paragraph{必要知識} 標準重力9.80665m/s\textsuperscript{2}、kgf=\allowbreak{}9.80665N、1cm\textsuperscript{2}=10\textasciicircum{}-4m\textsuperscript{2}、Pa=\allowbreak{}N/m\textsuperscript{2}、atm/bar/mmHg/psiの定義と慣用係数。

\paragraph{考え方・正答を選ぶ理由（査読待ち初稿）} 3が2桁遠い。\(1\mathrm{kgf\,cm^{-2}}=\frac{9.80665\mathrm{N}}{10^{-4}\mathrm{m^2}}=9.80665\times10^4\mathrm{Pa}\)である。原典の10\textsuperscript{2}と係数を比較する。kgfとkg、面積のcm\textsuperscript{2}とcmを混同しない。

\Needspace{4\baselineskip}\paragraph{選択肢の検討（査読待ち初稿）}\mbox{}\par

\noindent\textbf{1：適切} \(1\mathrm{atm}=1.01325\times10^5\mathrm{Pa}\)。標準気圧の定義係数である。実際の当地大気圧が常にこの値とはしない。\par

\noindent\textbf{2：適切} \(1\mathrm{bar}=10^5\mathrm{Pa}\)。barは10\textasciicircum{}5Paの単位定義である。\par

\noindent\textbf{3：不適切} \(1\mathrm{kgf\,cm^{-2}}=9.80665\times10^2\mathrm{Pa}\)。kgfは標準重力9.80665m/s\textsuperscript{2}で1kgに働く力9.80665N。cm\textsuperscript{2}=10\textasciicircum{}-4m\textsuperscript{2}で割ると9.80665e4Pa、原典は100分の1である。\par

\noindent\textbf{4：適切} \(1\mathrm{mmHg}=133.322\mathrm{Pa}\)。通常の水銀柱換算は133.322Pa程度、液柱の実測は密度・温度・重力条件を持つ。\par

\noindent\textbf{5：適切} \(1\mathrm{psi}=6.894757\times10^3\mathrm{Pa}\)。psiはlbf/in\textsuperscript{2}で約6894.757Pa。NIST通常換算と一致する有効桁である。\par

\paragraph{再利用する考え方} 定義の基準日を確定し、物理量の関係から指数と係数を導く。量子現象では運動量・エネルギー保存と観測対象を区別する。

\paragraph{共通理論} \hyperref[theory:units:conversion]{第\ref*{theory:units:conversion}節}, \hyperref[theory:units:derived]{第\ref*{theory:units:derived}節}

\par 参考：\cite{si_nist811}。

\par 難易度 2（校準初稿）、優先度未評価。

\section{SI基本単位の定義・改定}

\subsection{平成28年度 問6}
\label{q:h28-reg-applied-06}

\noindent\textbf{原典：} sources/H28/17\_applied\_science.pdf、PDF物理2頁。

\textbf{公式正答：} 4（single）。

\paragraph{問題の要旨} 改定前のSI基本単位の定義から、温度基準の取り違えを見つける。

\Needspace{4\baselineskip}\paragraph{確認した条件・記述（独自要約）}\mbox{}\par

\noindent\textbf{方向：} 最も不適切なもの。\par

\noindent\textbf{対象：} SIの定義・換算。定義の時点と指数を全肢画像で照合。\par

\paragraph{問われる核心} 時点に応じたSI基本単位の基準

\paragraph{必要知識} 旧Kの水三重点273.16K、旧kg原器・mol炭素12基準、2019-05-20改定、現行h/e/k/NA固定。

\paragraph{考え方・正答を選ぶ理由（査読待ち初稿）} 4が水の凝固点と三重点を取り違える。2019改定前のKは同位体組成を規定した水三重点の熱力学温度の1/273.16。通常圧凝固点とは異なる。現在はボルツマン定数kを固定する。出題年は2016であるため、kg原器と12C基準のmolを現在の知識だけで誤答扱いしない。定義は測定手順と区別する。

\Needspace{4\baselineskip}\paragraph{選択肢の検討（査読待ち初稿）}\mbox{}\par

\noindent\textbf{1：適切} 秒：Cs133超微細構造遷移に対応する放射の周期。。出題当時はCs133基底状態の二超微細準位間の放射9192631770周期を秒とした。現行も同じ周波数定数を固定する。\par

\noindent\textbf{2：適切} m：定めた時間に真空を進む光の距離。。真空光速cを299792458m/sに固定し、秒との関係からmを定義する。\par

\noindent\textbf{3：適切} kg：国際キログラム原器の質量。。2016当時のkgは原器定義で適切。2019-05-20後はh固定となり原器は定義ではない。\par

\noindent\textbf{4：不適切} K：水の凝固点を基準。。2019改定前のKは同位体組成を規定した水三重点の熱力学温度の1/273.16。通常圧凝固点とは異なる。現在はボルツマン定数kを固定する。\par

\noindent\textbf{5：適切} mol：.012kgの12C中の原子数と同数の要素粒子を含む物質量。。2016当時のmolはこの炭素12基準で適切。現在はNAを固定し、12gの12Cを定義文として使わない。\par

\paragraph{再利用する考え方} 定義の基準日を確定し、物理量の関係から指数と係数を導く。量子現象では運動量・エネルギー保存と観測対象を区別する。

\paragraph{共通理論} \hyperref[theory:units:history]{第\ref*{theory:units:history}節}, \hyperref[theory:units:definition]{第\ref*{theory:units:definition}節}, \hyperref[theory:units:temperature]{第\ref*{theory:units:temperature}節}

\par 参考：\cite{si_nist811,si_kelvin2007,si_cgpm2018}。

\paragraph{関連問題} \hyperref[q:h29-reg-applied-05]{H29-REG-APPLIED-05}

\par 難易度 2（校準初稿）、優先度未評価。

\subsection{平成29年度 問5}
\label{q:h29-reg-applied-05}

\noindent\textbf{原典：} sources/H29/17\_applied\_science.pdf、PDF物理2頁。

\textbf{公式正答：} 4（single）。

\paragraph{問題の要旨} 改定前の温度基準と電位・角度の単位関係を比較する。

\Needspace{4\baselineskip}\paragraph{確認した条件・記述（独自要約）}\mbox{}\par

\noindent\textbf{方向：} 最も不適切なもの。\par

\noindent\textbf{対象：} SIの定義・換算。定義の時点と指数を全肢画像で照合。\par

\paragraph{問われる核心} 温度の定義時点と電位のエネルギー基準

\paragraph{必要知識} 旧KとCelsiusオフセット、Q=\allowbreak{}It、V=仕事/Q、電場F/Q、rad=弧長/半径。

\paragraph{考え方・正答を選ぶ理由（査読待ち初稿）} 4が電位差と電場を取り違える。\(V=E_{\rm work}/Q\)よりV=\allowbreak{}J/C=\allowbreak{}Nm/Cである。水三重点の1・2は出題時の旧K定義として読む。平面角は長さ比でrad=\allowbreak{}1の次元を持つが角の意味を保持する。

\Needspace{4\baselineskip}\paragraph{選択肢の検討（査読待ち初稿）}\mbox{}\par

\noindent\textbf{1：適切} Kは水三重点で定義。。2019改定前のKは同位体組成を規定した水三重点の熱力学温度の1/273.16。通常圧凝固点とは異なる。現在はボルツマン定数kを固定する。\par

\noindent\textbf{2：適切} 水三重点は.01℃。。2017の旧定義では三重点273.16KとCelsiusの273.15Kオフセットから.01℃。定義水組成条件を持つ。\par

\noindent\textbf{3：適切} 1C=\allowbreak{}As。。電荷Q=\allowbreak{}ItよりC=\allowbreak{}Asである。\par

\noindent\textbf{4：不適切} 1V=\allowbreak{}1N C\textasciicircum{}-1。。N/Cは電場の次元である。VはJ/C=\allowbreak{}Nm/Cなので長さが一つ不足する。\par

\noindent\textbf{5：適切} 平面角はradで無次元。。平面角は弧長/半径で次元1、単位radを用いる。radは角度の意味を残す記号であり、角度でない任意の比と区別する。\par

\paragraph{再利用する考え方} 定義の基準日を確定し、物理量の関係から指数と係数を導く。量子現象では運動量・エネルギー保存と観測対象を区別する。

\paragraph{共通理論} \hyperref[theory:units:history]{第\ref*{theory:units:history}節}, \hyperref[theory:units:temperature]{第\ref*{theory:units:temperature}節}, \hyperref[theory:units:electrical]{第\ref*{theory:units:electrical}節}, \hyperref[theory:units:angle]{第\ref*{theory:units:angle}節}

\par 参考：\cite{si_kelvin2007,si_nist811,si_cgpm2018}。

\paragraph{関連問題} \hyperref[q:h28-reg-applied-06]{H28-REG-APPLIED-06}, \hyperref[q:r01-ret-applied-01]{R01-RET-APPLIED-01}

\par 難易度 2（校準初稿）、優先度未評価。

\subsection{令和元年度再試験 問1}
\label{q:r01-ret-applied-01}

\noindent\textbf{原典：} sources/R01\_retest/17\_applied\_science.pdf、PDF物理1頁。

\textbf{公式正答：} 1（single）。

\par\textbf{検討未確定：} 改定後を明記する原典の2は三重点.01℃を無条件に断定する。熱力学温度Tは改定後不確かさを持ち、ITS-90のT90と同じ厳密定義値ではない。提示桁の近似としては妥当、厳密さが未指定である。公式・初回1は維持し原器誤りと別に疑義を保持。

\paragraph{問題の要旨} 2019年5月改定後のSIを旧原器定義と区別する。

\Needspace{4\baselineskip}\paragraph{確認した条件・記述（独自要約）}\mbox{}\par

\noindent\textbf{方向：} 最も不適切なもの。\par

\noindent\textbf{対象：} SIの定義・換算。定義の時点と指数を全肢画像で照合。\par

\paragraph{問われる核心} 定義定数と三重点の測定値の違い

\paragraph{必要知識} 新kgはh固定、秒と光速m、C=\allowbreak{}As、Celsiusオフセット273.15K、水三重点の定義から実測への変化、rad。

\paragraph{考え方・正答を選ぶ理由（査読待ち初稿）} 意図された誤りは1であり、2019年5月後はh固定である。2の.01℃は提示精度での約値としては整合するが、以前の厳密定義値ではない。公式1と初回1を保持し、2の精度・温度尺度の省略は別の疑義にする。

\Needspace{4\baselineskip}\paragraph{選択肢の検討（査読待ち初稿）}\mbox{}\par

\noindent\textbf{1：不適切} kgは国際キログラム原器の質量で定義。。原器定義は2019-05-20に廃止。プランク定数h=\allowbreak{}6.62607015e-34Jsを固定する。Js=\allowbreak{}kg m\textsuperscript{2}/sなのでmとsの既定義と結びkgを定義する。重力定数Gを固定しない。\par

\noindent\textbf{2：近似として適切・厳密さ疑義} 水三重点は.01℃。。最良推定は273.16Kすなわち.01℃であるが、新K定義では実測値と不確かさを持つ。ITS-90のT90固定点を真の熱力学温度の定義と混同しない。\par

\noindent\textbf{3：適切} 1C=\allowbreak{}1As。。電荷Q=\allowbreak{}ItよりC=\allowbreak{}Asである。\par

\noindent\textbf{4：適切} 1mは真空光が1/299792458秒に進む距離。。真空光速cを299792458m/sに固定し、秒との関係からmを定義する。\par

\noindent\textbf{5：適切} 平面角はradで無次元。。平面角は弧長/半径で次元1、単位radを用いる。radは角度の意味を残す記号であり、角度でない任意の比と区別する。\par

\paragraph{再利用する考え方} 定義の基準日を確定し、物理量の関係から指数と係数を導く。量子現象では運動量・エネルギー保存と観測対象を区別する。

\paragraph{共通理論} \hyperref[theory:units:history]{第\ref*{theory:units:history}節}, \hyperref[theory:units:definition]{第\ref*{theory:units:definition}節}, \hyperref[theory:units:temperature]{第\ref*{theory:units:temperature}節}, \hyperref[theory:units:angle]{第\ref*{theory:units:angle}節}

\par 参考：\cite{si_cgpm2018,si_kelvin_realization,si_brochure9}。

\paragraph{関連問題} \hyperref[q:h29-reg-applied-05]{H29-REG-APPLIED-05}

\par 難易度 2（校準初稿）、優先度未評価。

\subsection{令和2年度 問9}
\label{q:r02-reg-applied-09}

\noindent\textbf{原典：} sources/R02/17\_applied\_science.pdf、PDF物理3頁。

\textbf{公式正答：} 2（single）。

\paragraph{問題の要旨} 現行SIの基本単位と定義定数の対応を判定する。

\Needspace{4\baselineskip}\paragraph{確認した条件・記述（独自要約）}\mbox{}\par

\noindent\textbf{方向：} 最も不適切なもの。\par

\noindent\textbf{対象：} SIの定義・換算。定義の時点と指数を全肢画像で照合。\par

\paragraph{問われる核心} 基本単位の定数による定義

\paragraph{必要知識} 2019改定のh/e/k/NA、Cs周波数、c、単色540THz視感効果度、定義と実現の区別。

\paragraph{考え方・正答を選ぶ理由（査読待ち初稿）} 2のkgと重力定数Gの対応が誤り。プランク定数h=\allowbreak{}6.62607015e-34Jsを固定する。Js=\allowbreak{}kg m\textsuperscript{2}/sなのでmとsの既定義と結びkgを定義する。重力定数Gを固定しない。 定数一つが他単位から独立に一単位を定義するわけではない。既に固定された時間・長さ等の体系と結ぶ。

\Needspace{4\baselineskip}\paragraph{選択肢の検討（査読待ち初稿）}\mbox{}\par

\noindent\textbf{1：適切} mは真空光速で定義。。真空光速cを299792458m/sに固定し、秒との関係からmを定義する。\par

\noindent\textbf{2：不適切} kgは重力定数で定義。。プランク定数h=\allowbreak{}6.62607015e-34Jsを固定する。Js=\allowbreak{}kg m\textsuperscript{2}/sなのでmとsの既定義と結びkgを定義する。重力定数Gを固定しない。原典の重力Gとの対応が誤り。\par

\noindent\textbf{3：適切} Aは電気素量で定義。。素電荷e=\allowbreak{}1.602176634e-19Cを固定する。C=\allowbreak{}Asなので秒と電荷の関係でAを定義する。\par

\noindent\textbf{4：適切} Kはボルツマン定数で定義。。k=\allowbreak{}1.380649e-23J/Kを固定し、エネルギー尺度と温度尺度を結ぶ。\par

\noindent\textbf{5：適切} molはアボガドロ定数で定義。。NA=\allowbreak{}6.02214076e23mol\textasciicircum{}-1を固定し、1molが含む指定要素粒子数を定める。粒子の種類を明示する。\par

\paragraph{再利用する考え方} 定義の基準日を確定し、物理量の関係から指数と係数を導く。量子現象では運動量・エネルギー保存と観測対象を区別する。

\paragraph{共通理論} \hyperref[theory:units:definition]{第\ref*{theory:units:definition}節}, \hyperref[theory:units:history]{第\ref*{theory:units:history}節}, \hyperref[theory:units:photometry]{第\ref*{theory:units:photometry}節}

\par 参考：\cite{si_cgpm2018,si_brochure9}。

\paragraph{関連問題} \hyperref[q:r07-reg-applied-09]{R07-REG-APPLIED-09}

\par 難易度 2（校準初稿）、優先度未評価。

\subsection{令和3年度 問9}
\label{q:r03-reg-applied-09}

\noindent\textbf{原典：} sources/R03/17\_applied\_science.pdf、PDF物理4頁。

\textbf{公式正答：} 1（single）。

\paragraph{問題の要旨} 現行SIの基本単位と定義定数の対応を判定する。

\Needspace{4\baselineskip}\paragraph{確認した条件・記述（独自要約）}\mbox{}\par

\noindent\textbf{方向：} 最も不適切なもの。\par

\noindent\textbf{対象：} SIの定義・換算。定義の時点と指数を全肢画像で照合。\par

\paragraph{問われる核心} 基本単位の定数による定義

\paragraph{必要知識} 2019改定のh/e/k/NA、Cs周波数、c、単色540THz視感効果度、定義と実現の区別。

\paragraph{考え方・正答を選ぶ理由（査読待ち初稿）} 黒体を使う1が誤り。540THzの単色放射の視感効果度Kcd=\allowbreak{}683lm/Wを固定する。光度は方向当たりの量で、現在は黒体温度を定義基準としない。 定数一つが他単位から独立に一単位を定義するわけではない。既に固定された時間・長さ等の体系と結ぶ。

\Needspace{4\baselineskip}\paragraph{選択肢の検討（査読待ち初稿）}\mbox{}\par

\noindent\textbf{1：不適切} cdは黒体放射で定義。。540THzの単色放射の視感効果度Kcd=\allowbreak{}683lm/Wを固定する。光度は方向当たりの量で、現在は黒体温度を定義基準としない。黒体基準は1979以前の歴史的定義に関係する。\par

\noindent\textbf{2：適切} mは真空光速で定義。。真空光速cを299792458m/sに固定し、秒との関係からmを定義する。\par

\noindent\textbf{3：適切} kgはプランク定数で定義。。プランク定数h=\allowbreak{}6.62607015e-34Jsを固定する。Js=\allowbreak{}kg m\textsuperscript{2}/sなのでmとsの既定義と結びkgを定義する。重力定数Gを固定しない。\par

\noindent\textbf{4：適切} 秒はCs超微細遷移周波数で定義。。Cs133の摂動のない基底状態超微細遷移周波数9192631770Hzを固定する。\par

\noindent\textbf{5：適切} molはアボガドロ定数で定義。。NA=\allowbreak{}6.02214076e23mol\textasciicircum{}-1を固定し、1molが含む指定要素粒子数を定める。粒子の種類を明示する。\par

\paragraph{再利用する考え方} 定義の基準日を確定し、物理量の関係から指数と係数を導く。量子現象では運動量・エネルギー保存と観測対象を区別する。

\paragraph{共通理論} \hyperref[theory:units:definition]{第\ref*{theory:units:definition}節}, \hyperref[theory:units:history]{第\ref*{theory:units:history}節}, \hyperref[theory:units:photometry]{第\ref*{theory:units:photometry}節}

\par 参考：\cite{si_cgpm2018,si_brochure9}。

\paragraph{関連問題} \hyperref[q:r07-reg-applied-09]{R07-REG-APPLIED-09}

\par 難易度 2（校準初稿）、優先度未評価。

\subsection{令和7年度 問9}
\label{q:r07-reg-applied-09}

\noindent\textbf{原典：} sources/R07/17\_applied\_science.pdf、PDF物理3頁。

\textbf{公式正答：} 3（single）。

\paragraph{問題の要旨} 現行SIの基本単位と定義定数の対応を判定する。

\Needspace{4\baselineskip}\paragraph{確認した条件・記述（独自要約）}\mbox{}\par

\noindent\textbf{方向：} 最も不適切なもの。\par

\noindent\textbf{対象：} SIの定義・換算。定義の時点と指数を全肢画像で照合。\par

\paragraph{問われる核心} 基本単位の定数による定義

\paragraph{必要知識} 2019改定のh/e/k/NA、Cs周波数、c、単色540THz視感効果度、定義と実現の区別。

\paragraph{考え方・正答を選ぶ理由（査読待ち初稿）} 3のkgと重力定数Gの対応が誤り。プランク定数h=\allowbreak{}6.62607015e-34Jsを固定する。Js=\allowbreak{}kg m\textsuperscript{2}/sなのでmとsの既定義と結びkgを定義する。重力定数Gを固定しない。 定数一つが他単位から独立に一単位を定義するわけではない。既に固定された時間・長さ等の体系と結ぶ。

\Needspace{4\baselineskip}\paragraph{選択肢の検討（査読待ち初稿）}\mbox{}\par

\noindent\textbf{1：適切} mは真空光速で定義。。真空光速cを299792458m/sに固定し、秒との関係からmを定義する。\par

\noindent\textbf{2：適切} Aは素電荷で定義。。素電荷e=\allowbreak{}1.602176634e-19Cを固定する。C=\allowbreak{}Asなので秒と電荷の関係でAを定義する。\par

\noindent\textbf{3：不適切} kgは重力定数で定義。。プランク定数h=\allowbreak{}6.62607015e-34Jsを固定する。Js=\allowbreak{}kg m\textsuperscript{2}/sなのでmとsの既定義と結びkgを定義する。重力定数Gを固定しない。原典の重力Gとの対応が誤り。\par

\noindent\textbf{4：適切} molはアボガドロ定数で定義。。NA=\allowbreak{}6.02214076e23mol\textasciicircum{}-1を固定し、1molが含む指定要素粒子数を定める。粒子の種類を明示する。\par

\noindent\textbf{5：適切} cdは単色放射で定義。。540THzの単色放射の視感効果度Kcd=\allowbreak{}683lm/Wを固定する。光度は方向当たりの量で、現在は黒体温度を定義基準としない。\par

\paragraph{再利用する考え方} 定義の基準日を確定し、物理量の関係から指数と係数を導く。量子現象では運動量・エネルギー保存と観測対象を区別する。

\paragraph{共通理論} \hyperref[theory:units:definition]{第\ref*{theory:units:definition}節}, \hyperref[theory:units:history]{第\ref*{theory:units:history}節}, \hyperref[theory:units:photometry]{第\ref*{theory:units:photometry}節}

\par 参考：\cite{si_cgpm2018,si_brochure9}。

\paragraph{関連問題} \hyperref[q:r02-reg-applied-09]{R02-REG-APPLIED-09}, \hyperref[q:r03-reg-applied-09]{R03-REG-APPLIED-09}

\par 難易度 2（校準初稿）、優先度未評価。
